\documentclass[11pt]{article}
\usepackage[margin=1in]{geometry}
\usepackage[T1]{fontenc}
\usepackage{lmodern}
\usepackage{hyperref}
\usepackage{enumitem}
\usepackage{booktabs}
\usepackage{array}
\usepackage{titlesec}

\hypersetup{colorlinks=true, linkcolor=blue, urlcolor=blue}
\titleformat{\section}{\large\bfseries}{\thesection}{0.6em}{}
\setlist{itemsep=2pt, topsep=4pt}

\title{Project 1: Agentic App Research and Project 2 Proposal}
\author{Anthony Bayate\\Kean University \,--\, Web Client-Side Programming}
\date{October 2026}

\begin{document}
\maketitle

\section{Introduction}
An AI agent differs from a plain chatbot in one way: it does not just answer, it
\emph{acts}. It is given a goal, decides which tools to call (search, read a file,
run code, call an API), looks at what came back, and keeps going until the goal is
met or it needs a human. Below I look at five existing AI-native apps that use at
least one agent, and for each I describe how the agent is built in, why the app
needs it, and what it actually does. After that I propose my own web app for
Project 2.

\section{Five Existing Agentic Apps}

\subsection{Cursor (agent mode)}
\textbf{What it does.} Cursor is a code editor built on VS Code. In agent mode you
describe a task (``add a login page and wire it to the API'') and the agent plans
the change, reads the relevant files, edits several files, runs terminal commands,
and reads the errors that come back.

\textbf{How it is integrated.} The agent lives inside the editor, so its tools are
editor features: file search, file read/write, a terminal, and the linter output.
The model calls these tools in a loop, and every edit is shown as a diff that the
developer accepts or rejects.

\textbf{Why it is needed.} A normal autocomplete suggestion only sees the line you
are typing. Real features touch many files, and the developer would otherwise do the
searching and copy-pasting by hand. The agent removes that busywork and uses the
compiler and test errors to correct itself.

\subsection{GitHub Copilot coding agent}
\textbf{What it does.} You assign a GitHub issue to Copilot. It spins up a
sandboxed environment, works on the repository, and opens a pull request with its
changes for you to review.

\textbf{How it is integrated.} It is wired into GitHub itself: issues are the input,
Actions-style runners are the sandbox, and the pull request is the output. The
review step is where the human stays in control.

\textbf{Why it is needed.} Many issues are small and well described (fix a bug, add
a test, update a dependency). Handing those off lets a developer spend time on the
design work instead. Running asynchronously in a sandbox also means a mistake cannot
touch the main branch.

\subsection{Perplexity (research and Deep Research)}
\textbf{What it does.} Perplexity answers questions with cited sources. In its deep
research mode, the agent breaks a question into sub-questions, runs many searches,
reads the pages, and writes a structured report with links.

\textbf{How it is integrated.} Search is the core tool, not an add-on. The model
decides what to search next based on what it has already read, then synthesizes the
result and attaches citations.

\textbf{Why it is needed.} A single search or a single model call cannot cover a
broad question, and a model without search can be out of date or make things up.
The agent loop gives fresher answers and lets the user check the sources.

\subsection{Devin (Cognition)}
\textbf{What it does.} Devin is marketed as an autonomous software engineering
agent. It gets a task in a chat or from a ticket, then works in its own environment
with a shell, a code editor, and a browser. It can set up a project, write code, run
tests, and report back.

\textbf{How it is integrated.} Instead of living in an existing tool, Devin is the
workspace: the user watches a timeline of its commands and can step in with
feedback.

\textbf{Why it is needed.} It is aimed at longer tasks that take many steps, such
as migrations or repetitive refactors, where the main cost is a person's time rather
than difficulty. The tradeoff is that long autonomous runs need good visibility and
review so mistakes are caught.

\subsection{Claude Code}
\textbf{What it does.} Claude Code is a terminal-based coding agent. It reads a
codebase, edits files, runs commands and tests, and uses git. It can also call
outside tools through MCP servers and hand pieces of work to subagents.

\textbf{How it is integrated.} It runs in the developer's own terminal and
project folder with a permission system: reading is easy, but running commands or
editing files can require approval. Project instructions live in a plain
\texttt{CLAUDE.md} file.

\textbf{Why it is needed.} A lot of development happens outside an editor window,
in scripts, builds, and git. An agent with shell access can do that whole loop of
change, run, read the error, and fix, while the permission model keeps the user in
charge of risky actions.

\subsection{What these apps have in common}
\begin{itemize}
  \item \textbf{A goal, not a prompt.} The user states an outcome and the agent
        works out the steps.
  \item \textbf{Tools.} Each one gets its power from tools (files, shell, search,
        APIs) rather than from the model alone.
  \item \textbf{Feedback loop.} Errors, test results, and search results go back
        into the model so it can correct course.
  \item \textbf{A human checkpoint.} Diffs, pull requests, and permission prompts
        keep a person in the loop for anything that matters.
\end{itemize}

\section{My Proposal for Project 2: StudyPilot}

\subsection{The problem}
College students juggle syllabi, assignment deadlines, and exams across several
classes, and the information is scattered over PDFs, LMS pages, and emails. Most
students build a plan by hand once and then abandon it when the semester gets busy.
A chatbot can give generic study advice, but it does not know \emph{my} deadlines
and cannot adjust when my week changes.

\subsection{The idea}
StudyPilot is a web app where a student uploads their course syllabi (or pastes
assignment lists). An AI agent turns that into a living weekly study plan, quizzes
the student on the material, and re-plans when something changes (a missed session,
a new due date, a low quiz score).

\subsection{What the agent does}
The agent receives a goal such as ``get me ready for Friday's midterm'' and chooses
among its tools:
\begin{enumerate}
  \item \textbf{Parse syllabus:} extract courses, due dates, exam dates, and weights
        into structured data.
  \item \textbf{Plan:} place study blocks on a calendar based on how much time the
        student has and how heavy each deadline is.
  \item \textbf{Quiz:} generate practice questions from the student's own notes and
        grade the answers.
  \item \textbf{Re-plan:} if quiz scores are weak or a block is skipped, move time
        toward the weak topic and update the calendar.
  \item \textbf{Ask:} when information is missing or a choice is ambiguous (for
        example, two exams the same day), ask the student instead of guessing.
\end{enumerate}

\subsection{Why an agent and not just a chatbot}
The task is a loop. The plan depends on quiz results, and quiz results change the
plan. A single prompt cannot do that because it has no way to read the calendar,
write to it, or remember what happened last week. Tool use plus state is the
reason this needs an agent.

\subsection{Proposed architecture}
\begin{center}
\begin{tabular}{@{}p{0.2\linewidth}p{0.7\linewidth}@{}}
\toprule
\textbf{Layer} & \textbf{Plan} \\
\midrule
Front end & HTML, CSS, and JavaScript (possibly React): upload page, weekly
calendar view, quiz view, chat panel showing the agent's actions. \\
Back end & Node.js/Express server. The API key stays on the server and is never
shipped to the browser. \\
Agent & An LLM API with function (tool) calling; a loop that runs tools until the
goal is met or the user is needed. \\
Tools & \texttt{parse\_syllabus}, \texttt{get\_calendar}, \texttt{add\_study\_block},
\texttt{generate\_quiz}, \texttt{grade\_answer}, \texttt{ask\_user}. \\
Storage & A database for courses, deadlines, plans, and quiz history. \\
\bottomrule
\end{tabular}
\end{center}

\subsection{Safety and limits}
\begin{itemize}
  \item The agent proposes calendar changes and the student confirms them before
        anything is saved.
  \item Uploaded files are used only for that student's plan.
  \item Every tool call is logged and visible in the UI, so the student can see what
        the agent did and why.
  \item Limits on steps per run to avoid runaway loops and unexpected API costs.
\end{itemize}

\subsection{Milestones}
\begin{enumerate}
  \item Basic web UI and server with one LLM call (syllabus to JSON).
  \item Add tool calling and the planning loop.
  \item Add quizzes, grading, and re-planning.
  \item Polish the UI, add confirmations and logging, and test with real syllabi.
\end{enumerate}

\subsection{How I will know it works}
I will test it with real syllabi from my own classes and check three things: the
extracted dates match the documents, the plan has no conflicts, and the re-plan
actually shifts time after a poor quiz score.

\section{Conclusion}
The five apps above all follow the same pattern: a goal, a set of tools, a feedback
loop, and a human checkpoint. StudyPilot applies that pattern to a problem I deal
with every semester, and it is small enough to build in the time available.

\section*{References}
\small
\textbf{Blogs and documentation}
\begin{enumerate}[label={[\arabic*]}]
  \item Anthropic, ``Building effective agents.''
        \url{https://www.anthropic.com/research/building-effective-agents}
  \item Cursor documentation (Agent mode). \url{https://cursor.com/docs}
  \item GitHub Docs, ``About Copilot coding agent.''
        \url{https://docs.github.com/en/copilot}
  \item Perplexity, ``Introducing Perplexity Deep Research.''
        \url{https://www.perplexity.ai/hub/blog/introducing-perplexity-deep-research}
  \item Cognition, ``Introducing Devin.'' \url{https://cognition.ai/blog/introducing-devin}
  \item Anthropic, Claude Code documentation.
        \url{https://docs.anthropic.com/en/docs/claude-code/overview}
  \item OpenAI, ``Function calling'' guide.
        \url{https://platform.openai.com/docs/guides/function-calling}
\end{enumerate}
\textbf{GitHub repositories}
\begin{enumerate}[label={[\arabic*]}, resume]
  \item Anthropic, Claude Code. \url{https://github.com/anthropics/claude-code}
  \item OpenAI Agents SDK (Python). \url{https://github.com/openai/openai-agents-python}
\end{enumerate}
\textbf{Videos and courses}
\begin{enumerate}[label={[\arabic*]}, resume]
  \item Kaggle \& Google, 5-Day AI Agents Intensive course podcast (YouTube).
        \url{https://www.youtube.com/watch?v=cbzmr7vt4XA}
  \item Kaggle \& Google, 5-Day AI Agents Intensive / Vibe Coding course.
        \url{https://www.kaggle.com/competitions/5-day-ai-agents-intensive-vibecoding-course-with-google}
\end{enumerate}

\end{document}
